\subsection{NOTATION}

\begin{enumerate}
\def\labelenumi{(\alph{enumi})}
\tightlist
\item
  As a general rule the law of a commutative monoid is written additively. It is then a convention that -x denotes the negative of x. The notation \(x + (-y)\) is abbreviated to x - y and similarly
\end{enumerate}

\[
x + y - z, \quad x - y - z, \quad x - y + z - t, \quad \text { etc. } \dots
\]

represent respectively

\[
x + y + (- z), \quad x + (- y) + (- z), \quad x + (- y) + z + (- t), \quad \text { etc. } \dots
\]

For \(n \in \mathbf{Z}\) the notation \(\top^n x\) is replaced by \(nx\) . Formulae (9) to (12) then become

\[
(m + n). x = m. x + n. x\tag{13}
\]

\[
0. x = 0\tag{14}
\]

\[
1. x = x\tag{15}
\]

\[
m. (n. x) = (m n). x\tag{16}
\]

where m and n belong to N or even to Z if x admits a negative. Also in the latter case the relation \((-1)\) . x = -x holds. We also note the formula

\[
n. (x + y) = n. x + n. y.\tag{17}
\]

\begin{enumerate}
\def\labelenumi{(\alph{enumi})}
\setcounter{enumi}{1}
\tightlist
\item
  Let \(\mathbf{E}\) be a monoid written multiplicatively. For \(n \in \mathbf{Z}\) the notation \(\frac{n}{T} x\) is replaced by \(x^n\) . We have the relations
\end{enumerate}

\[
\begin{array}{r l} x ^ {m + n} & = x ^ {m}. x ^ {n} \\ x ^ {0} & = 1 \\ x ^ {1} & = x \\ (x ^ {m}) ^ {n} & = x ^ {m n} \end{array}
\]

and also \((xy)^n = x^n y^n\) if \(x\) and \(y\) commute.

When \(x\) has an inverse, this is precisely \(x^{-1}\) . The notation \(\frac{1}{x}\) is also used instead of \(x^{-1}\) . Finally, when the monoid \(E\) is commutative, \(\frac{x}{y}\) or \(x / y\) is also used for \(xy^{-1}\) .

